%%%%%%%%%%%%%%%%%%%%%%%%%%%%%%%%%%%%%%%%
\chapter{Conclusiones}
%%%%%%%%%%%%%%%%%%%%%%%%%%%%%%%%%%%%%%%%

La presente revisión sistemática analizó 84 estudios empíricos publicados entre 2020 y 2025 sobre técnicas de \textit{fine-tuning} para \textit{Large Language Models}, seleccionados mediante el protocolo PRISMA desde una búsqueda inicial de 140 registros en Scopus. Los hallazgos se organizan en correspondencia con los tres objetivos específicos planteados.

En relación con el primer objetivo, sistematizar la literatura mediante PRISMA extrayendo características técnicas, métricas y resultados por dominio, el corpus presenta una distribución heterogénea entre cuatro dominios de aplicación: Salud y Biomedicina (30 estudios), Lenguaje y Humanidades (21), Ciencias Exactas y Computación (20) e Industria y Negocios (13). LoRA y sus variantes, incluyendo QLoRA, emergen como la categoría dominante con presencia en el 79,8\% de los estudios (67 de 84). La calidad metodológica del corpus fue aceptable en conjunto, con un \textit{score} promedio de 2,61 sobre 3,0. La extracción de datos permitió documentar técnicas, modelos base, \textit{datasets}, métricas y resultados de forma estructurada y trazable para cada uno de los 84 estudios incluidos, cumpliendo los requisitos de reproducibilidad del protocolo PRISMA.

En relación con el segundo objetivo, analizar el rendimiento por dominio, los \textit{trade-offs} y los patrones de aplicabilidad, el análisis revela que el \textit{fine-tuning} paramétrico eficiente produce mejoras consistentes sobre el modelo base sin adaptar en todos los dominios evaluados, aunque con magnitudes variables según el contexto. Los estudios con métricas comparables muestran que las mejoras más consistentes se concentran en tareas de baja cobertura lingüística y dominios altamente especializados, donde el ajuste fino sobre datos del dominio produce ganancias sustanciales independientemente del tamaño del modelo. Los patrones de aplicabilidad difieren entre dominios: en Salud y Biomedicina predomina la preferencia por modelos abiertos motivada por restricciones de privacidad; en Lenguaje y Humanidades destaca la diversificación de variantes de LoRA orientadas a tareas específicas; en Ciencias Exactas y Computación LoRA se consolida como referencia establecida e integrada con aprendizaje federado; y en Industria y Negocios la configuración dominante combina LoRA con RLHF para tareas de generación de alto valor.

En relación con el tercer objetivo, sintetizar criterios de selección según contexto operativo, limitaciones y direcciones futuras, la evidencia acumulada permite formular tres criterios basados en evidencia empírica. Cuando existen restricciones de privacidad o se opera sobre datos sensibles, QLoRA sobre modelos abiertos es la configuración más respaldada por la literatura, con documentación explícita en estudios clínicos y biomédicos. Cuando el dominio de aplicación es altamente especializado o la lengua tiene baja cobertura en los datos de preentrenamiento, LoRA con \textit{datasets} específicos del dominio produce las mejoras más sustanciales independientemente del tamaño del modelo base. Cuando el objetivo es la calidad y alineación del texto generado en aplicaciones de producción, la combinación LoRA + RLHF es la configuración predominante en la literatura reciente, especialmente en dominios industriales y legales.

En respuesta a la pregunta de investigación, las técnicas de \textit{fine-tuning} para \textit{Large Language Models} desarrolladas entre 2020 y 2025 convergen hacia un modelo de adaptación paramétrica eficiente con LoRA como componente central, complementado con mecanismos adicionales según el contexto operativo. La selección de técnica no depende exclusivamente del rendimiento absoluto sino de factores contextuales como la privacidad de los datos, la disponibilidad de recursos lingüísticos del dominio y los requisitos de alineación de las respuestas generadas. Este patrón, consistente a través de dominios y tareas, sugiere que la comunidad investigadora ha convergido hacia un conjunto acotado de configuraciones efectivas durante el período analizado.

La presente revisión presenta cuatro limitaciones principales que deben considerarse al interpretar sus hallazgos. La primera es la restricción a una única base de datos: la búsqueda se realizó exclusivamente en Scopus, lo que puede haber excluido estudios relevantes indexados únicamente en IEEE Xplore, ACM Digital Library o disponibles como preprints en arXiv. La segunda es la heterogeneidad de \textit{baselines}: los estudios incluidos emplean referencias de comparación de naturaleza diversa, lo que impide una síntesis cuantitativa unificada y limita la comparación directa de magnitudes de mejora entre estudios y dominios. La tercera es la ausencia de \textit{baseline} explícito en el 64,3\% de los estudios incluidos (54 de 84), lo que concentra las comparaciones en el subconjunto con reporte completo. La cuarta es el sesgo de acceso abierto: el filtro aplicado para garantizar la reproducibilidad del proceso de extracción puede haber introducido un sesgo de selección hacia ciertos tipos de publicaciones o dominios con mayor tradición de publicación abierta.

Los hallazgos de esta revisión sugieren tres direcciones prioritarias para la investigación futura. La primera es la estandarización de protocolos de evaluación y reporte: la heterogeneidad de métricas, \textit{baselines} y configuraciones experimentales documentada en este corpus dificulta la acumulación de evidencia comparativa, y el desarrollo de estándares análogos a los existentes en otras áreas de \textit{machine learning} permitiría síntesis cuantitativas más robustas. La segunda es la investigación en escenarios de recursos limitados: los hallazgos más consistentes del corpus corresponden a adaptaciones en lenguas de baja cobertura y dominios con datos escasos, y la investigación futura debería profundizar en el tamaño mínimo de \textit{dataset} necesario para obtener mejoras significativas y en la sensibilidad a la calidad de los datos de entrenamiento. La tercera es la evaluación comparativa de técnicas de alineación: RLHF y DPO aparecen en el corpus principalmente en dominios aplicados pero con escasa evidencia cuantitativa comparativa entre ellas, y estudios que las comparen sistemáticamente con métricas y \textit{baselines} compartidos contribuirían a orientar la selección de técnica en aplicaciones de producción.